\begin{abstract}
Many multi-agent LLM papers claim that the improvements they observe follow from the
transformations implied by the role prompts (verifier, critic, devil's advocate) their agents
use, often without checking whether these prompts implement them. We propose executable
manipulation checks, drawing on the manipulation check paradigm in experimental psychology:
predeclared behavioral checks, compared to a length-matched neutral baseline.

We test six role prompts in three versions each: an explicit implementation, a bare role-name,
and a length-matched neutral prompt. Of six diagnostic checks on 200 examples each, one meets
the pre-set statistical criteria, one reverses, and three show no difference from neutral; we
performed a sixth check that was non-diagnostic after audit.

In terms of raw accuracy on the 1,000 scored HaluEval items, the three versions of the council
vary within a single percentage point. None shows a significant effect on leaderboard score,
according to our pre-registered test, and all score lower than a single call to the model.
However, bare role-name prompts do show major reductions in refusal rates, suggesting that role
prompting may affect refusal behavior, even if it doesn't implement its specified manipulation.

We show silent failures in some of the coalition attribution analysis used in this literature,
on simulated records. Retraining ablations may produce a Shapley value of zero that still meets
efficiency. Masking approximations may also not be ranking-preserving. A post-collection audit
further shows that manipulation checks themselves need degeneracy diagnostics: of six
pre-registered checks, five included specific guards; two were defeatable, three didn't vary on
observed support, and the one check that passed had none. Manipulation checks should survive
executions.
\end{abstract}
